In this section, we discuss the exploration of \pred\ (\gt) in the linear and non-linear new arms bandit setting discussed in \Cref{sec:new-actions}.  Recall that we consider the linear bandit setting of horizon $n=50$, $\Mpr = 200000$, $\Mts = 200$, $A=20$, and $d=5$. 
%
Here during data collection and during collecting the test data, we randomly select one new action from $\R^d$ for each task $m$. 
%
So the number of invariant actions is $|\Anc| = 19$.
% 
%


\textbf{Outcomes:} We first discuss the main outcomes of our analysis of exploration in the low-data regime:

\begin{tcolorbox}
    \customfinding The \pred\ (\gt) is robust to changes when the number of in-variant actions is large. \pred\ (\gt) performance drops as shared structure breaks down.
\end{tcolorbox}

% \noindent
% \textbf{(1)} The \pred\ (\gt) is robust to changes when the number of changing arms is small. It samples the optimal action in each task more than \dptg. 

% \textbf{(2)} When the number of in-variant actions is small then all algorithms \pred, \predt, and \ad\ sample the optimal action very less. This shows that the latent shared structure across the tasks is very important.



We first show in \Cref{fig:train-dist-new-1} the training distribution of the optimal actions. For each bar, the frequency indicates the number of tasks where the action (shown in the x-axis) is the optimal action.

Then in \Cref{fig:analysis-epxploration-new-1} we show how the sampling distribution of \dptg, \pred\, and \predt\ change in the first $10$ and last $10$ rounds for all the tasks where action $17$ is optimal. We plot this graph the same way as discussed in \Cref{sec:linear}.
%
% Again to plot this graph we first sum over the individual pulls of the action taken by each algorithm over the first $10$ and last $10$ rounds. Then we average these counts over all test tasks where action $5$ is optimal. 
%
From the figure \Cref{fig:analysis-epxploration-new-1} we see that \pred (\gt) consistently pulls the action $17$ more than \dptg. It also explores other optimal actions like $\{1,2,3,8,9,15\}$ but discards them quickly in favor of the optimal action $17$ in these tasks. 

Finally, we plot the feasible action set considered by \dptg, \pred, and \predt\ in \Cref{fig:analysis-exploration-time-new-1}. To plot this graph again we consider the test tasks where the optimal action is $17$. Then we count the number of distinct actions that are taken from round $t$ up until horizon $n$. Finally we average this over all the considered tasks where the optimal action is $17$. We call this the candidate action set considered by the algorithm. From the \Cref{fig:analysis-exploration-time-new-1} we see that \predt\ explores more than \pred\ in this setting.

% \dptg\ explores the least and gets stuck with few actions quickly (by round $10$). \predt\ explores more than \pred\ in this setting. We see that \dptg\ explores most but focuses more on sub-optimal actions (see \Cref{fig:analysis-epxploration-new-1}).


We also show how the prediction error of the optimal action by \pred\ compared to \linucb\ in this $1$ new arm linear bandit setting. In \Cref{fig:short-horizon-action-dist-new-1} we first show how the $20$ actions are distributed in the $\Mts=200$ test tasks. In \Cref{fig:short-horizon-action-dist-new-1} for each bar, the frequency indicates the number of tasks where the action (shown in the x-axis) is the optimal action. Then in \Cref{fig:short-horizon-action-error-new-1} we show the prediction error of \pred\ (\gt) for each task $m\in[\Mts]$. The prediction error is calculated the same way as stated in \Cref{sec:new-actions} 
%
From the \Cref{fig:short-horizon-action-error-new-1} we see that for most actions the prediction error of \pred\ (\gt) is closer to \linucb\ showing that the introduction of $1$ new action does not alter the prediction error much.
%
% actions $\{2,3,5,6,7,10\}$ the prediction error of \pred\ is either close or smaller than \linucb. 
% This constitutes nearly $70\%$ of the $200$ test tasks. 
Note that \linucb\ estimates the empirical mean directly from the test task, whereas \pred\ has a strong prior based on the training data. Therefore we see that \pred\ is able to estimate the reward of the optimal action quite well from the training dataset $\Dpr$.
%

% \begin{figure}[!hbt]
% \centering
% \begin{tabular}{ccc}
% \subfigure[Train Optimal Action Distribution]{\includegraphics[scale =0.27]{img/Neurips/new_arms/linear/final/linear_small/optimal_action_train.png}\label{fig:train-dist-new-1}} &
% %%%%%%%%%%
% \subfigure[Distribution of action sampling in all tasks where action $17$ is optimal]{\includegraphics[scale =0.27]{img/Neurips/new_arms/linear/final/linear_small/analysis_pull.png} \label{fig:analysis-epxploration-new-1}} &
% %%%%%%%%%%%%%%%%%%%%%%%
% \subfigure[Candidate Action Set in Time averaged over all tasks where action $17$ is optimal]{\hspace*{-1.5em}\includegraphics[scale =0.27]{img/Neurips/new_arms/linear/final/linear_small/analysis_time_1.png}\label{fig:analysis-exploration-time-new-1}} 
% \end{tabular}
% \caption{Exploration Analysis of \pred (\gt) in linear 1 new arm setting}
% \label{fig:exploration-analysis-time-new-1}
% \end{figure}
% %
% %
% %
% \begin{figure}[!hbt]
% \centering
% \vspace*{-1em}
% \begin{tabular}{cc}
% %%%%%%%%%%%%%%%%%%%%%%%
% \hspace*{-1.5em}\subfigure[Test action distribution]{\includegraphics[scale =0.27]{img/Neurips/new_arms/linear/final/linear_small/optimal_action_test.png} \label{fig:short-horizon-action-dist-new-1}}&
% % %%%%%%%%%%
% \hspace*{-1.5em}\subfigure[Test Prediction Error]{\includegraphics[scale =0.27]{img/Neurips/new_arms/linear/final/linear_small/analysis_total.png} \label{fig:short-horizon-action-error-new-1}}
% \end{tabular}
% \vspace*{-1em}
% \caption{Prediction error of \pred (\gt) in linear 1 new arm setting}
% \label{fig:expt-short-horizon-new-1}
% \vspace{-0.7em}
% \end{figure}


\begin{figure}[!hbt]
\centering
\begin{subfigure}[b]{0.32\textwidth}
    \includegraphics[scale=0.27]{img/Neurips/new_arms/linear/final/linear_small/optimal_action_train.png}
    \caption{Train Optimal Action Distribution}
    \label{fig:train-dist-new-1}
\end{subfigure}%
\begin{subfigure}[b]{0.32\textwidth}
    \includegraphics[scale=0.27]{img/Neurips/new_arms/linear/final/linear_small/analysis_pull.png}
    \caption{Distribution of action sampling in all tasks where action $17$ is optimal}
    \label{fig:analysis-epxploration-new-1}
\end{subfigure}%
\begin{subfigure}[b]{0.32\textwidth}
    \includegraphics[scale=0.27]{img/Neurips/new_arms/linear/final/linear_small/analysis_time_1.png}
    \caption{Candidate Action Set in Time averaged over all tasks where action $17$ is optimal}
    \label{fig:analysis-exploration-time-new-1}
\end{subfigure}
\caption{Exploration Analysis of \pred (\gt) in linear 1 new arm setting}
\label{fig:exploration-analysis-time-new-1}
\end{figure}

\begin{figure}[!hbt]
\centering
\vspace*{-1em}
\begin{subfigure}[b]{0.48\textwidth}
    \includegraphics[scale=0.27]{img/Neurips/new_arms/linear/final/linear_small/optimal_action_test.png}
    \caption{Test action distribution}
    \label{fig:short-horizon-action-dist-new-1}
\end{subfigure}%
\begin{subfigure}[b]{0.48\textwidth}
    \includegraphics[scale=0.27]{img/Neurips/new_arms/linear/final/linear_small/analysis_total.png}
    \caption{Test Prediction Error}
    \label{fig:short-horizon-action-error-new-1}
\end{subfigure}
\vspace*{-1em}
\caption{Prediction error of \pred (\gt) in linear 1 new arm setting}
\label{fig:expt-short-horizon-new-1}
\vspace{-0.7em}
\end{figure}

We now consider the setting where the number of invariant actions is $|\Anc| = 15$.
% 
%
We again show in \Cref{fig:train-dist-new-5} the training distribution of the optimal actions. For each bar, the frequency indicates the number of tasks where the action (shown in the x-axis) is the optimal action. Then in \Cref{fig:analysis-epxploration-new-5} we show how the sampling distribution of \dptg, \pred\, and \predt\ change in the first $10$ and last $10$ rounds for all the tasks where action $17$ is optimal. We plot this graph the same way as discussed in \Cref{sec:linear}.
%
% Again to plot this graph we first sum over the individual pulls of the action taken by each algorithm over the first $10$ and last $10$ rounds. Then we average these counts over all test tasks where action $5$ is optimal. 
%
From the figure \Cref{fig:analysis-epxploration-new-5} we see that none of the algorithms \pred, \predt, \dptg\ consistently pulls the action $17$ more than other actions. This shows that the common underlying actions across the tasks matter for learning the epxloration.
% It also explores other optimal actions like $\{1,2,3,8,9, 15\}$ but discards them quickly in favor of the optimal action $17$ in these tasks. 

Finally, we plot the feasible action set considered by \dptg, \pred, and \predt\ in \Cref{fig:analysis-exploration-time-new-5}. To plot this graph again we consider the test tasks where the optimal action is $17$. We build the candidate set the same way as before. From the \Cref{fig:analysis-exploration-time-new-5} we see that none of the three algorithms \dptg, \pred, \predt, is able to sample the optimal action $17$ sufficiently high number of times.

% \dptg\ explores the least and gets stuck with few actions quickly (by round $10$). \predt\ explores more than \pred\ in this setting. We see that \dptg\ explores most but focuses more on sub-optimal actions (see \Cref{fig:analysis-epxploration-new-1}).


We also show how the prediction error of the optimal action by \pred\ compared to \linucb\ in this $1$ new arm linear bandit setting. In \Cref{fig:short-horizon-action-dist-new-5} we first show how the $20$ actions are distributed in the $\Mts=200$ test tasks. In \Cref{fig:short-horizon-action-dist-new-5} for each bar, the frequency indicates the number of tasks where the action (shown in the x-axis) is the optimal action. Then in \Cref{fig:short-horizon-action-error-new-5} we show the prediction error of \pred\ (\gt) for each task $m\in[\Mts]$. The prediction error is calculated the same way as stated in \Cref{sec:new-actions}. 
%
From the \Cref{fig:short-horizon-action-error-new-5} we see that for most actions the prediction error is higher than \linucb\ showing that the introduction of $5$ new actions (and thereby decreasing the invariant action set) significantly alters the prediction error.
%
% actions $\{2,3,5,6,7,10\}$ the prediction error of \pred\ is either close or smaller than \linucb. 
% This constitutes nearly $70\%$ of the $200$ test tasks. 
% Note that \linucb\ estimates the empirical mean directly from the test task, whereas \pred\ has a strong prior based on the training data. Therefore we see that \pred\ is able to estimate the reward of the optimal action quite well from the training dataset $\Dpr$.


% \begin{figure}[!hbt]
% \centering
% \begin{tabular}{ccc}
% \subfigure[Train Optimal Action Distribution]{\includegraphics[scale = 0.27]{img/Neurips/new_arms/linear/linear_large/optimal_action_train.png}\label{fig:train-dist-new-5}} &
% %%%%%%%%%%
% \subfigure[Distribution of action sampling in all tasks where action $17$ is optimal]{\includegraphics[scale = 0.27]{img/Neurips/new_arms/linear/linear_large/analysis_pull.png} \label{fig:analysis-epxploration-new-5}} &
% %%%%%%%%%%%%%%%%%%%%%%%
% \subfigure[Candidate Action Set in Time averaged all tasks where action $17$ is optimal]{\includegraphics[scale = 0.27]{img/Neurips/new_arms/linear/linear_large/analysis_time_1.png}\label{fig:analysis-exploration-time-new-5}} 
% \end{tabular}
% \caption{Exploration Analysis of \pred (\gt) in linear 5 new arm setting}
% \label{fig:exploration-analysis-time-new-5}
% \end{figure}
% %
% %
% %
% \begin{figure}[!hbt]
% \centering
% \vspace*{-1em}
% \begin{tabular}{cc}
% %%%%%%%%%%%%%%%%%%%%%%%
% \hspace*{-1.5em}\subfigure[Test action distribution]{\includegraphics[scale =0.27]{img/Neurips/new_arms/linear/linear_large/optimal_action_test.png} \label{fig:short-horizon-action-dist-new-5}}&
% % %%%%%%%%%%
% \hspace*{-1.5em}\subfigure[Test Prediction Error]{\includegraphics[scale =0.27]{img/Neurips/new_arms/linear/linear_large/analysis_total.png} \label{fig:short-horizon-action-error-new-5}}
% \end{tabular}
% \vspace*{-1em}
% \caption{Prediction error of \pred (\gt) in linear 1 new arm setting}
% \label{fig:expt-short-horizon-new-5}
% \vspace{-0.7em}
% \end{figure}


\begin{figure}[!hbt]
\centering
\begin{subfigure}[b]{0.32\textwidth}
   \includegraphics[scale=0.27]{img/Neurips/new_arms/linear/linear_large/optimal_action_train.png}
   \caption{Train Optimal Action Distribution}
   \label{fig:train-dist-new-5}
\end{subfigure}%
\begin{subfigure}[b]{0.32\textwidth}
   \includegraphics[scale=0.27]{img/Neurips/new_arms/linear/linear_large/analysis_pull.png}
   \caption{Distribution of action sampling in all tasks where action $17$ is optimal}
   \label{fig:analysis-epxploration-new-5}
\end{subfigure}%
\begin{subfigure}[b]{0.32\textwidth}
   \includegraphics[scale=0.27]{img/Neurips/new_arms/linear/linear_large/analysis_time_1.png}
   \caption{Candidate Action Set in Time averaged all tasks where action $17$ is optimal}
   \label{fig:analysis-exploration-time-new-5}
\end{subfigure}
\caption{Exploration Analysis of \pred (\gt) in linear 5 new arm setting}
\label{fig:exploration-analysis-time-new-5}
\end{figure}

\begin{figure}[!hbt]
\centering
\vspace*{-1em}
\begin{subfigure}[b]{0.48\textwidth}
   \includegraphics[scale=0.27]{img/Neurips/new_arms/linear/linear_large/optimal_action_test.png}
   \caption{Test action distribution}
   \label{fig:short-horizon-action-dist-new-5}
\end{subfigure}%
\begin{subfigure}[b]{0.48\textwidth}
   \includegraphics[scale=0.27]{img/Neurips/new_arms/linear/linear_large/analysis_total.png}
   \caption{Test Prediction Error}
   \label{fig:short-horizon-action-error-new-5}
\end{subfigure}
\vspace*{-1em}
\caption{Prediction error of \pred (\gt) in linear 1 new arm setting}
\label{fig:expt-short-horizon-new-5}
\vspace{-0.7em}
\end{figure}